\newpage
\section{Conclusion et perspectives}
\subsection{Résumé du travail réalisé}
Au final, on a réussi à développer un jeu d'échecs chinois fonctionnel avec les éléments suivants :
\begin{itemize}
\item Un plateau 11×11 avec colonnes neutres et rivière, qui respecte toutes les spécifications ;
\item Les règles de déplacement de toutes les pièces, avec la gestion correcte des colonnes neutres ;
\item Un moteur de jeu qui valide chaque coup, détecte l'échec, l'échec et mat, et la règle du face à face ;
\item Un mode deux joueurs et un mode contre le bot avec trois niveaux de difficulté ;
\item Une interface graphique Swing avec historique des coups et affichage des captures ;
\item Un système de logs avec Log4j et des tests unitaires avec JUnit.
\end{itemize}
Ce projet nous a vraiment aidés à comprendre l'importance de bien séparer les responsabilités dans un logiciel, et à travailler en équipe sur un projet Java de taille réelle.
\subsection{Améliorations possibles du projet}
Si on avait eu plus de temps, on aurait voulu :
\begin{itemize}
\item Améliorer le bot avec un algorithme Minimax et l'élagage alpha-bêta pour le rendre plus fort ;
\item Ajouter une fonctionnalité de sauvegarde et de chargement de partie ;
\item Ajouter la rotation automatique du plateau en mode deux joueurs pour que chaque joueur voie son côté ;
\item Ajouter des animations pour les déplacements de pièces ;
\item Permettre de jouer en réseau contre un autre joueur.
\end{itemize}